%% kapa: exam=2 round=25 year=2014 period=1 subject=appraisal_practice subject_ko=감정평가실무 seq=4 number=4 points=20
%% source: Q-Net 감정평가사 기출문제 (cst00309) artlSeq=5109209 file=extracted/1008/5109209/2014년 제25회 감정평가사 2차시험 기출문제/1교시(감정평가실무).hwp
%% conversion: kordoc 4.24.51 -> tools/gen2.py (auto); review=proofread-sonnet; flags=none
\begin{longquestion}{【문제 4】}
다음 물음에 답하시오. (20점) \par
\medskip\noindent 물음 1) A국에서 과거 약 20년간 상업용 부동산(A,B,C)과 주식(D)의 연간평균수익률 추이를 조사한 후, 다음의 표와 같이 정리하였다. \par
(단위 : \%) \par
{\small\centering
\begin{tabular}{|c|c|c|c|c|c|}
\hline
\makecell{구분\newline{}(자료기호)} & \makecell{작성 기초자료의 성격} & \makecell{기하평균\newline{}수익률} & \makecell{산술평균\newline{}수익률} & \makecell{표준\newline{}편차} & \makecell{시계열\newline{}상관계수} \\
\hline
\makecell{CREF (A)} & \makecell{매년의 감정평가액 집계} & \makecell{10.8} & \makecell{10.9} & \makecell{2.6} & \makecell{0.43} \\
\hline
\makecell{REITs (B)} & \makecell{리츠의 수익률 집계} & \makecell{14.2} & \makecell{15.7} & \makecell{15.4} & \makecell{0.11} \\
\hline
\makecell{C\&S (C)} & \makecell{실거래가격의 통계처리} & \makecell{8.5} & \makecell{8.6} & \makecell{3.0} & \makecell{0.17} \\
\hline
\makecell{S\&P 500 (D)} & \makecell{주요주식의 거래가격} & \makecell{12.3} & \makecell{13.5} & \makecell{16.7} & \makecell{-0.10} \\
\hline
\end{tabular}\par}

이 자료의 해석과 관련하여, (1) 기하평균수익률과 산술평균수익률이 상이할 때 무엇을 채택하는 것이 합리적인지, (2) 상기 자료 B와 자료 D는 표준편차가 유사하고, 시계열상관계수도 낮은 경향을 보인 반면 자료A는 표준편차가 가장 낮고, 시계열 상관계수가 가장 높은 특징을 보이는 이유를 기초자료의 성격과 관련하여 약술하시오. (10점) \par
\medskip\noindent 물음 2) 완전소유권의 시장가치는 임대권가치와 임차권가치의 합이라 할 때, 연간 시장임대료(순임료)는 12,000,000원, 연간 계약임대료(순임료)는 9,000,000원, 계약기간 10년, 계약기간 만료 시 본건 부동산의 완전소유권 시장가치는 120,000,000원이고 계약기간 중 시장가치의 변동은 없는 것으로 예상되는 경우 다음 물음에 답하시오. (10점) \par
\medskip\noindent (1) 임대권(賃貸權) 수익률이 9.00\%라고 할 경우 본건의 내재된 임차권(賃借權) 수익률은 얼마인가? (5점) \par
\medskip\noindent (2) 어떤 경우에「완전소유권 시장가치=임대권가치+임차권가치」의 등식이 성립하지 않는가? (5점) \par
\end{longquestion}
